% ID: SEL-D0083-OA-E03
% ESTADO: cotejada. Segunda pasada visual y borrador compilado revisados.
% Fuente DOC: Selectividad/sel_2012_matematicas_aplicadas/5 MATEMÁT. APLIC. CC. SS/EXAMEN 3-ANDALUCÍA.doc
% SHA-256 original: c8688a19515dfeb348e8f8817d3f4a6c0a967723c270f2ec42e7f365d001e555
% Página convertida: 1; opción A; ejercicio 3.

Un pescador tiene tres tipos de carnada de las que sólo una es adecuada para pescar salmón. Si utiliza la carnada correcta la probabilidad de que pesque un salmón es $1/3$, mientras que si usa una de las inadecuadas esa probabilidad se reduce a $1/5$.
\begin{enumerate}[label=\alph*)]
\item \textbf{(1.25 puntos)} Si elige aleatoriamente la carnada, \textquestiondown{}cuál es la probabilidad de que pesque un salmón?
\item \textbf{(1.25 puntos)} Si ha pescado un salmón, \textquestiondown{}cuál es la probabilidad de que lo haya hecho con la carnada adecuada?
\end{enumerate}
